% !TeX program = xelatex
% =============================================================
% 模板 06：ACM SIGCONF 会议论文（acmart 官方文档类）
% 编译： scripts\build.ps1 en-acm
% 版式： 双栏，9 pt，官方 acmart.cls（TeX Live 自带）
% 说明： 使用 nonacm 选项关闭版权块，投稿时删掉 nonacm 并补全 \acmConference 等信息
% =============================================================
\documentclass[sigconf,nonacm]{acmart}

\usepackage{booktabs}
\usepackage{amsmath}

% 关闭 acmart 默认的页眉页脚信息，得到干净的预览版
\settopmatter{printfolios=true,printacmref=false}

\begin{document}

% ------------------------- 标题区（改这里） -------------------------
\title{Placeholder Title: A Sample ACM SIGCONF Paper for Template Testing}

\author{San Zhang}
\affiliation{%
  \institution{School of Computer Science, XX University}
  \city{Beijing}
  \country{China}}
\email{zhangsan@example.edu.cn}

\author{Si Li}
\affiliation{%
  \institution{School of Software, XX University}
  \city{Beijing}
  \country{China}}
\email{lisi@example.edu.cn}

\renewcommand{\shortauthors}{Zhang and Li}

\begin{abstract}
This is the placeholder abstract for an ACM SIGCONF paper. It shows the acmart
title block, abstract, CCS concepts, a numbered equation, a table, and a BibTeX
reference list. Replace the text with a 150--250 word abstract of your own work.
\end{abstract}

\begin{CCSXML}
<ccs2012>
   <concept>
       <concept_id>10010147.10010178</concept_id>
       <concept_desc>Computing methodologies~Machine learning</concept_desc>
       <concept_significance>300</concept_significance>
   </concept>
</ccs2012>
\end{CCSXML}

\ccsdesc[300]{Computing methodologies~Machine learning}
\keywords{placeholder, template, acmart, sigconf}

\maketitle

\section{Introduction}
This section is a placeholder introduction. State the problem, summarize prior
work~\cite{vaswani2017attention}, and end with a short list of contributions.

\section{Approach}
Given an input sequence $\mathbf{x}=(x_1,\dots,x_n)$, the attention output is
\begin{equation}
  \mathrm{Attn}(\mathbf{x})=\operatorname{softmax}\!\left(\frac{\mathbf{Q}\mathbf{K}^{\top}}{\sqrt{d_k}}\right)\mathbf{V},
  \label{eq:attn}
\end{equation}
where $\mathbf{Q},\mathbf{K},\mathbf{V}$ are the query, key and value matrices and
$d_k$ is the key dimension. Equation~\eqref{eq:attn} is the equation placeholder.

\section{Experiments}
Table~\ref{tab:setup} lists the placeholder experimental configuration.

\begin{table}[h]
  \caption{Placeholder Experimental Configuration}
  \label{tab:setup}
  \centering
  \begin{tabular}{lc}
    \toprule
    Item & Value \\
    \midrule
    Dataset      & Placeholder-1M \\
    Optimizer    & Adam           \\
    Learning rate& $3\times10^{-4}$ \\
    Batch size   & 64             \\
    \bottomrule
  \end{tabular}
\end{table}

\section{Results}
The placeholder system improves accuracy by 2.1 points over the baseline, which
follows the same direction as the results reported in~\cite{devlin2019bert}.

\section{Conclusion}
This template compiles with XeLaTeX and BibTeX using the official acmart class.
Replace all placeholder content before submission.

\bibliographystyle{ACM-Reference-Format}
\bibliography{refs}

\end{document}
